Hospital Information Systems (HIS) are the core infrastructure of the modern public health system. Accordingly, the development of a suitable educational HIS for academic training in human medicine and medical informatics is essential. To create a fully functional and realistic educational HIS, the project investigates the generation of structured, longitudinal synthetic clinical data alongside the integration of continuous sensor data streams from medical wearables.

The ability to proficiently navigate and use operational HIS is a fundamental competency in modern medical education \citep{BICHELFINDLAY2023104908}. As healthcare becomes increasingly digitized, students must not only understand traditional clinical workflows but also learn to manage heterogeneous data formats, use various technical interfaces, and evaluate the growing role of patient-generated health data in the treatment context \citep{healthcare13030329}. Practical, hands-on experience with realistic, evolving patient histories and continuous physiological monitoring is therefore important to effectively prepare future physicians for the complexities of modern, data-driven patient care \citep{Mohan2016-qn}.

Despite this educational necessity, traditional training environments face significant structural and legal limitations. The widespread use of authentic electronic health records (EHRs) is strictly restricted by data privacy regulations such as the General Data Protection Regulation (GDPR) and the Health Insurance Portability and Accountability Act (HIPAA) \citep{Chen2021-ji}. Consequently, academic HIS often rely on unpopulated databases or static synthetic datasets \citep{choi2021}. These static approaches represent merely a cross-section of a specific point in time, failing to adequately reflect the dynamic, longitudinal nature of clinical workflows or the continuous evolution of clinical parameters required to analyze disease progressions \citep{Majumder2017-tg}.

To bridge this pedagogical gap, this project aims to develop a methodology for providing dynamic data for a dedicated educational HIS within the Heidelberg \enquote{studiKIS} project. The primary objective is to synthetically generate structured clinical data that continuously evolves over time, enabling students to observe and analyze realistic, longitudinal treatment trajectories. Furthermore, the project integrates real-world sensor data to teach students how to process diverse data formats and technical interoperability standards, such as HL7 FHIR. By engineering this dynamic data pipeline and systematically evaluating its clinical plausibility, the project yields a privacy-compliant, highly realistic simulation environment that bridges the gap between theoretical knowledge and practical clinical application.